\scanpage{311}
\mspage{157}

L'ens. $\unsure{\widetilde U'|U'}$ des pts $Q\in U$ tels que $G_Q^+\ne(1)$ est stable par l'action de $\mathcal G$, et s'identifie (avec cette action) à \unsure{l'ens.} des relèvements \inserted{\unsure{maximaux}} mod $\pi$ de ss-groupes (cycliques) $\ne1$ de $G_{\unsure{Q}}^+$. Quand on connaît, pour un relèvement partiel $r$ d'un $G_Q$ dans $\mathcal E$ quel est son centralisateur dans $\mathcal E$ i.e. le relèvement correspondant de $\unsure{\mathcal G_k|\mathcal G_l}$, \unsure{bien} de $\unsure{\overline{\mathrm m}|=}$ pour les conjugués de $r$ par n'importe quel $\unsure{\text{él}^{\mathrm{mt}}}$ $g$ (non seulement de $\pi$, mais $\unsure{\overline{\mathrm m}|=}$ de $\mathcal E$), par simple conjugaison.

Donc les cas à déterminer correspondent pratiquement aux orbites de $\mathcal G$ dans $U'$.

On peut s'intéresser à décrire $\mathcal G$ ou $\Gamma$ en tant que ss-groupe de $\mathrm{Aut\,ext}_{\unsure{\mathrm{lac}}}(\pi)=\Gamma'$ \inserted{donnant lieu à l'ext. $\mathcal E'$ de $\Gamma'$ par $\pi$ (donc $\mathcal E\subset\mathcal E'$).} Mais pour tt relèvement $r$ d'un $G_Q$, considérons $\mathrm{Centr}_{\mathcal E'}(r)$, \deleted{\unsure{on}} : on aura $\mathrm{Centr}_{\mathcal E'}(r)\cap\pi=(1)$ i.e. \inserted{on trouve une section au dessus de} l'image de ce centralisateur dans $\Gamma'$, soit $\Gamma'(Q)$. Cette image ne dépend que de $Q$ i.e. de la classe de $\pi$-conjugaison de $r$ ou de $r(\unsure{G_Q})$, et est remplacée par un $G$-conjugué quand $Q$ est remplacé par un $G$-conjugué. Ceci

\mspage{310}

\scanpage{312}

dit, l'intersection \inserted{$\Gamma'^{\unsure{\natural|\sharp}}$} des $\Gamma'(Q)$, pour $Q\in U'$, est un ss-groupe de $\Gamma'$ qui contient l'intersection \deleted{\unsure{de}} $\mathcal G^{\unsure{\natural|\sharp}}$ des $\mathcal G_Q^{\deleted{\unsure{\natural}}}$, et le centralisateur de $G$ dans $\Gamma'^{\unsure{\natural|\sharp}}$ \unsure{contient} de $\unsure{\overline{\mathrm m}|=}$ l'intersection \deleted{de} $\Gamma^{\unsure{\natural|\sharp}}$ des $\Gamma_Q$, qui est un ss-groupe \inserted{invariant} d'indice fini de $\Gamma$. Et on peut \unsure{lors|alors} se proposer s'il est possible de construire au moins \unsure{le|un} ss-groupe fini $\Gamma^{\unsure{\natural|\sharp}}$ de $\Gamma$ comme $\mathrm{Centr}_{\Gamma'^{\unsure{\natural|\sharp}}}(G)$, et \unsure{d'en récupérer} peut être $\Gamma$ comme le normalisateur de $\Gamma^{\unsure{\natural|\sharp}}$ dans \deleted{le centralisateur} $\mathrm{Centr}_{\unsure{G|\mathcal G}}(\unsure{\Gamma|r})$.

Je m'intéresse plus particulièrement à la variante profinie de ceci, dans le cas où $U=\mathbb P^1_{\mathbb Q}\setminus\{0,1,\infty\}$, $G=\unsure{\mathfrak S_3|\mathcal S_3}$, $\Gamma={}$ groupe de Galois de $\overline{\mathbb Q}/\mathbb Q$, \deleted{\unsure{tandis}} \inserted{$\unsure{P|\rho}=\exp(2i\pi/6)$}.
\margin{clôture alg. \unsure{dans} $\mathbb C$}

Je n'ai pas vérifié que $\Gamma\longrightarrow\mathrm{Autext}_{\unsure{\mathrm{lac}}}(\widehat\pi)$ soit injectif, cela \unsure{n'empêche|m'empêche} de faire des calculs dans $\mathrm{Aut}_{\unsure{\mathrm{lac}}}(\widehat\pi)$, fussent-ils heuristiques pour le moment.

\mspage{311}

\scanpage{313}
\mspage{158}

Je \unsure{bute} \unsure{à|:} des ennuis de notations -- trop de groupes sont \unsure{\uline{désignés}} par la lettre $\Gamma$ (avec \unsure{évt.} des primes, indices, exposants \ldots). Il y a trois types de groupes qui interviennent dans mes réflexions :

a) les groupes de Teichmüller et ses variantes, qui jouent le rôle \unsure{des|de} groupes ``universels'' opérant (\inserted{éventuellement} mod isotopie) sur des surfaces, ou sur des groupes \deleted{\unsure{à lacets}} extérieurs à lacets. Ces groupes ont tendance à être infinis

b) Des groupes (le plus souvent finis) opérant sur des surfaces topologiques, ou sur des courbes algébriques géométriques (\unsure{sans en cela}, bouger le corps de base).

\margin{Les cas b) et c) se mélangent parfois (\unsure{dans} \unsure{$G$}, \unsure{groupes} \unsure{$\mathcal G$} ext. d'un groupe de Galois $\Gamma$ par un groupe \unsure{$G$} \deleted{\unsure{fini}} \inserted{Alg.} \unsure{dans le cas de} \unsure{$G$} \unsure{Géom.}}

c) Des groupes de Galois profinis (donc infinis), surtout des corps de type fini sur $\mathbb Q$, opérant ``arithmétiquement'' sur des surfaces et leurs $\widehat{\pi}_1$ géométriques.

C'est à cause des analogies profondes entre les cas b) et c), et leurs relations étroites avec le cas a), que j'avais été induit à adopter des notations communes, mais

\mspage{312}

\scanpage{314}

qui à la longue finissent par aboutir à des collisions. Il y a donc lieu de revoir les notations. Je vais réserver la lettre $G$ \inserted{et variantes} pour des actions géométriques (cas b)) de groupes, le plus souvent finis, la lettre $\Gamma$ \inserted{et variantes} pour des groupes de Galois, la lettre $\mathcal G$ pour des groupes mixtes.

\unsure{Quant|Quand} aux groupes ``universels'' de type Teichmüller, comme ceux notés $\Gamma_{g,\nu}$ précédemment, je vais plutôt les noter \inserted{$\mathcal T$,} $\mathcal T_{g,\nu}$ (initiale de ``Teichmüller'', alors que $\Gamma$, \inserted{$G$} \unsure{sont} l'initiale de Galois).

Le groupe de Galois \inserted{sur $\mathbb Q$} de la clôture algébrique $\overline{\mathbb Q}$ de $\mathbb Q$ dans $\mathbb C$ \deleted{\unsure{de $\overline{\mathbb Q}$ sur}} \unsure{mérite} une lettre spéciale, je le noterai $\boldsymbol{\Gamma}$. Le quotient \deleted{\unsure{par} $\widehat{\Gamma}_{g,\nu}$ de} $\operatorname{Norm}_{\widehat{\widehat{\mathcal T}}_{g,\nu}}(\widehat{\mathcal T}_{g,\nu})/\widehat{\mathcal T}_{g,\nu}$, qui s'apparente plus à un groupe de Galois qu'à un groupe de Teichmüller, sera noté $\boldsymbol{\Gamma}_{g,\nu}$ (lettre grasse !). Dans le

\mspage{313}

\scanpage{315}
\mspage{159}

cas $(g,\nu)=(0,3)$ qui m'occupe plus particulièrement, $\boldsymbol{\Gamma}_{0,3}$ s'identifie au \deleted{\unsure{normalisateur}} \inserted{centralisateur} de $G=\unsure{\mathfrak S_3|\mathcal S_3}$ \inserted{$=\mathcal T_{0,3}^{+}$} dans $\widehat{\widehat{\mathcal T}}_{0,3}$ -- il est contenu dans $\widehat{\widehat{\mathcal T}}_{03}^{!}$, \inserted{et peut-être égal. \unsure{?|à}}
\margin{$\widehat{\widehat{\mathcal T}}$ est ici produit $\frac{1}{2}$ direct de $\unsure{\mathfrak S_3|\mathcal S_3}=\mathcal T^+$ par $\mathcal T^{!}$}
on a un hom canonique
\[
\boldsymbol{\Gamma}\longrightarrow\boldsymbol{\Gamma}_{03}
\qquad (\text{plus généralement }\boldsymbol{\Gamma}\longrightarrow\boldsymbol{\Gamma}_{g,\nu})
\]
dont j'ignore pour l'instant s'il est injectif, et encore plus s'il est surjectif. Les réflexions qui précèdent suggèrent des conditions sur l'image, qui sont surtout intéressantes si on admet les relations
\[
\widehat{\pi}^{\unsure{\rho|\beta}}=\widehat{\pi}^{\sigma_0}=(1)
\]

On \inserted{a} \unsure{désigné\deleted{ra}} par $\mathcal E_{g\nu}$ l'extension canonique de $\mathcal T_{g,\nu}$ par $\pi_{g,\nu}$ (qui pour $(g,\nu)$ anabélien s'identifie à $\mathcal T'_{g,\nu+1}$).

L'opération extérieure d'un groupe $G$, $\Gamma$ ou $\mathcal G$ définit aussi une extension par $\pi$ (ou par $\widehat{\pi}$), qu'on a également \unsure{désigné|désignée} par la lettre $\mathcal E$ (initiale d'extension) -- il y a à nouveau collision de notations.

\mspage{314}

\scanpage{316}

Je vais prendre la lettre $\unsure{\mathfrak{S}|\mathfrak{G}}$ (qui fait pendant à $\mathcal{T}$) pour ces extensions dans les cas universels \unsure{à la} Teichmüller, (en écrivant $\unsure{\mathfrak{S}|\mathfrak{G}}_{g,\nu}$ au lieu des $\Gamma'_{g,\nu+1}$, puisque l'optique est différente\ldots), et en gardant la lettre $\mathcal{E}$ dans les cas \unsure{précédents}. Donc $\mathcal{E}$ a tendance à être une sous-extension d'une $\unsure{\mathfrak{S}|\mathfrak{G}}$.

\unsure{Admettons|Admettant} que $\boldsymbol{\Gamma}\longrightarrow\boldsymbol{\Gamma}_{0,3}$ est injectif, on aurait donc

\diagram{\begin{tikzcd}[ampersand replacement=\&,column sep=small,row sep=large]
1 \arrow[r,no head] \& \widehat{\pi}_{0,3} \arrow[r] \arrow[d,equal] \& \widehat{\widehat{\unsure{\mathfrak{S}|\mathfrak{G}}}}_{0,3} \arrow[r] \arrow[d] \& \widehat{\widehat{\mathcal{T}}}_{0,3} \arrow[r] \arrow[d] \& 1 \\
1 \arrow[r,no head] \& \widehat{\pi}_{0,3} \arrow[r] \arrow[d,equal] \& \widehat{\widehat{\unsure{\mathfrak{S}|\mathfrak{G}}}}'_{0,3} \arrow[r] \arrow[d] \& \underbrace{\unsure{\mathcal{T}^{+}_{0,3}}}_{\unsure{\mathfrak{S}_{3}}}\;\deleted{\unsure{\boldsymbol{\Gamma}}}\!\times\boldsymbol{\Gamma}_{0,3} \arrow[r] \arrow[d] \& 1 \\
1 \arrow[r] \& \widehat{\pi}_{03} \arrow[r] \arrow[d] \& \mathcal{E} \arrow[r] \arrow[d] \& \unsure{\mathfrak{S}_{3}}\times\boldsymbol{\Gamma} \arrow[r,no head] \arrow[d] \& 1 \\
1 \arrow[r] \& \pi_{0,3} \arrow[r] \& \unsure{\mathfrak{S}|\mathfrak{G}}_{03} \arrow[r,no head] \& \underbrace{\unsure{\mathfrak{S}_{3}}\times\unsure{\mathbb{Z}/2}}_{\mathcal{T}_{0,3}} \arrow[r] \& 1
\end{tikzcd}}

\margin{NB On note $\widehat{\widehat{\unsure{\mathfrak{S}|\mathfrak{G}}}}'_{g,\nu}$ le normalisateur de $\widehat{\unsure{\mathfrak{S}|\mathfrak{G}}}^{+}_{g,\nu}$ dans $\widehat{\widehat{\unsure{\mathfrak{S}|\mathfrak{G}}}}_{g,\nu}$, extension de $\boldsymbol{\Gamma}_{g,\nu}$ par $\widehat{\unsure{\mathfrak{S}|\mathfrak{G}}}^{+}_{g,\nu}$}

\deleted{\unsure{Soit}} $\gamma\in\boldsymbol{\Gamma}_{0,3}$, \deleted{\unsure{qui soit}}\inserted{pour qu'il soit} dans l'image de $\boldsymbol{\Gamma}$ \deleted{\unsure{ait|est} la propriété}\inserted{il faut} \unsure{qu'il|qu'ils} \unsure{admette|admettent} un relèvement $u$ qui commute à $\rho$, et un rel. qui commute à $\sigma_{0}$. \unsure{Ce|(ce} qui, \unsure{dès|des} que $\gamma\in\widehat{\widehat{\mathcal{T}}}_{0,3}$, implique

\mspage{315}

\scanpage{317}
\mspage{160}

déjà que $\unsure{\sigma|\gamma}$ dans $\widehat{\widehat{\mathcal{T}}}_{0,3}$ commute à $\mathcal{T}^{+}_{0,3}=\unsure{\mathfrak{S}_{3}}$, \deleted{donc} i.e. qu'il est dans $\boldsymbol{\Gamma}_{0,3}$). Il se pourrait que \unsure{tout} élément de $\boldsymbol{\Gamma}_{0,3}$ ait déjà cette propriété, donc que cette condition ne \unsure{fasse} pas de restriction véritable sur l'image de $\boldsymbol{\Gamma}$ dans $\boldsymbol{\Gamma}_{0,3}$.

\mspage{316}

\scanpage{318}
\mspage{161}

\uline{26.} Groupes de Teichmüller profinis.

\margin{\unsure{Discrétifications} et \unsure{prodiscrétifications|prédiscrétifications}.}

Soit $\pi$ un groupe profini à lacets \inserted{de type $g,\nu$}, $T$ le $\widehat{\mathbb Z}$-module inversible de ses orientations.

On suppose qu'on est dans le cas \unsure{anabélien}, et on admettra \unsure{qu'alors}
\[
\operatorname{Centre}(\pi)=\{1\}
\tag{1}
\]
plus généralement que le centralisateur dans $\pi$ de tout sous-groupe ouvert de $\pi$ est réduit à $\{1\}$. On aura donc encore une suite exacte canonique de groupes profinis
\[
1 \unsure{\longrightarrow|{-}} \pi \longrightarrow \operatorname{Aut}_{\mathrm{lac}}(\pi) \longrightarrow \operatorname{Autext}_{\mathrm{lac}}(\pi) \longrightarrow 1
\tag{\unsure{2}}
\]
On posera aussi \deleted{\unsure{$\operatorname{Autext}$}}
\[
\operatorname{Autext}_{\mathrm{lac}}(\pi)=\widehat{\widehat{\mathcal T}}(\pi)
\tag{3}
\]
en l'appelant le \uline{groupe de Teichmüller étendu} de $\pi$. On posera aussi $\widehat{\widehat{\unsure{\mathfrak S|\mathfrak G}}}(\pi)=\operatorname{Aut}_{\mathrm{lac}}(\pi)$, de sorte qu'on peut écrire (2) comme
\[
1 \longrightarrow \pi \longrightarrow \widehat{\widehat{\unsure{\mathfrak S|\mathfrak G}}}(\pi) \longrightarrow \widehat{\widehat{\mathcal T}}(\pi) \longrightarrow 1
\tag{2'}
\]

Appelons "base" de $\pi$ un \unsure{ens} d'éléments de $\pi$ $(x_i,y_i)_{1\leq i\leq g}$, $(\ell_j)_{1\leq j\leq\nu}$, \inserted{les $\ell_j$ \unsure{engend$^{\mathrm t}$} les diff. gr. à lacets} satisfaisant
\[
\deleted{\unsure{\prod_i}\,\unsure{[x_i,y_i]}\,\unsure{\cdots[x_g,y_g]\ell_1}}
\quad
\inserted{\ell_\nu\ell_{\nu-1}\cdots\ell_1[x_g,y_g]\cdots[x_1,y_1]=1}
\tag{\unsure{4}}
\]
et tels que ceci soit une relation génératrice.

\mspage{317}

\scanpage{319}

\margin{\unsure{ou encore}, on définit $\Pi_{g,\nu}$ comme le groupe discret \unsure{de|des} gén. les $x_i,y_i,\ell_j$ et de relation de \unsure{définition} (4)}

\unsure{Si} on choisit dans $\Pi_{g,\nu}$ une base (discrète) (\unsure{définition|Définition} correspondante), \unsure{d'où} une base de $\widehat{\Pi}_{g,\nu}$, on aura une bijection évidente
\[
\operatorname{Bases}(\pi) \xleftarrow{\sim} \operatorname{Isom}_{\mathrm{lac}}(\widehat{\Pi}_{g,\nu},\pi)
\tag{5}
\]

L'ensemble des bases de $\pi$ est un ensemble homogène \unsure{dans|sous} \deleted{\illegible{1}} $\widehat{\widehat{\unsure{\mathfrak S|\mathfrak G}}}(\pi)=\operatorname{Aut}_{\mathrm{lac}}(\pi)$. \unsure{Soit|Si} la \deleted{\unsure{$\Pi_0$ d}} base correspond à un iso $u:\widehat{\Pi}_{g,\nu}\longrightarrow\pi$, le groupe \inserted{$\Pi_0$} engendré par les $x_i,y_i,\ell_j$ n'est autre que $\deleted{\unsure{\Pi_0}}\inserted{u}(\Pi_{g,\nu})$. \deleted{Une des} \inserted{Les sous-groupes} de $\pi$ qui peuvent s'obtenir ainsi sont appelés les \uline{discrétifications} de $\pi$. Celles-ci forment un ensemble homogène sous $\widehat{\widehat{\unsure{\mathfrak S|\mathfrak G}}}=\operatorname{Aut}_{\mathrm{lac}}(\pi)$, \inserted{can.} isomorphe au quotient du $\widehat{\widehat{\unsure{\mathfrak S|\mathfrak G}}}_{g,\nu}$-\deleted{\unsure{torseur}}\inserted{ensemble} à droite $\operatorname{Isom}_{\mathrm{lac}}(\widehat{\Pi}_{g,\nu},\pi)$ par $\unsure{\mathfrak S|\mathfrak G}_{g,\nu}$ :
\[
\unsure{\operatorname{Discrt}|\operatorname{Discret}}(\pi) \xleftarrow{\sim} \operatorname{Isom}_{\mathrm{lac}}(\widehat{\Pi}_{g,\nu},\pi)/\unsure{\mathfrak S|\mathfrak G}_{g,\nu}
\tag{6}
\]
Les \unsure{autom} \inserted{(profinis à lacets)} de $\pi$ fixant \inserted{une discrétification} $\Pi_0$ s'identifient aux \unsure{autom} du groupe discret \inserted{à lacets} $\Pi_0$.
\[
\operatorname{Aut}_{\mathrm{lac}}(\pi,\Pi_0)\simeq\operatorname{Aut}_{\mathrm{lac}}(\Pi_0)
\tag{7}
\]

\mspage{318}

\scanpage{320}
\mspage{162}

La bijection (5) met sur l'ens des \unsure{bases|Bases} de $\pi$ une structure de bitorseur \unsure{sous|sur} $\widehat{\widehat{\unsure{\mathfrak S|\mathfrak G}}}(\pi),\widehat{\widehat{\unsure{\mathfrak S|\mathfrak G}}}_{g,\nu}$, et \unsure{les|la} \unsure{topologies|topologie} \unsure{correspondantes|correspondante} en fait un ens profini. L'ensemble des discrétifications de $\pi$, qui est un quotient de l'ens précédent, hérite d'une topologie quotient, qui n'est autre que la top. quotient du deuxième membre de (6). \deleted{Sauf dans les cas \unsure{exceptionnels} \unsure{$(g,\nu)=(0,3)$} \unsure{où} $\unsure{\mathfrak S|\mathfrak G}_{g,\nu}$} \inserted{(NB $\unsure{\mathfrak S|\mathfrak G}_{g,\nu}$ est \unsure{tjs} infini)} \inserted{Cet espace n'est} pas séparé, car le groupe $\unsure{\mathfrak S|\mathfrak G}_{g,\nu}\subset\widehat{\widehat{\unsure{\mathfrak S|\mathfrak G}}}_{g,\nu}$ n'est pas fermé. On désigne par $\unsure{\widetilde{\operatorname{Discrt}}'|\operatorname{Discrt}'}(\pi)$ l'espace topologique séparé associé, s'identifiant à $\operatorname{Isom}_{\mathrm{lac}}(\widehat{\Pi}_{g,\nu},\pi)/\unsure{\overline{\unsure{\mathfrak S|\mathfrak G}}}_{g,\nu}$, où $\overline{\unsure{\mathfrak S|\mathfrak G}}_{g,\nu}$ désigne l'adhérence de $\unsure{\mathfrak S|\mathfrak G}_{g,\nu}$ dans $\widehat{\widehat{\unsure{\mathfrak S|\mathfrak G}}}_{g,\nu}$. On a d'ailleurs un hom surjectif évident $\widehat{\unsure{\mathfrak S|\mathfrak G}}_{g,\nu}\longrightarrow\unsure{\widehat{\widehat{\unsure{\mathfrak S|\mathfrak G}}}|\overline{\unsure{\mathfrak S|\mathfrak G}}}_{g,\nu}$ dont l'image est $\overline{\unsure{\mathfrak S|\mathfrak G}}_{g,\nu}$, \unsure{nous admettrons} qu'il est injectif, et identifierons $\overline{\unsure{\mathfrak S|\mathfrak G}}_{g,\nu}$ à $\widehat{\unsure{\mathfrak S|\mathfrak G}}_{g,\nu}$. Ainsi
\[
\unsure{\operatorname{Discrt}|\operatorname{Discret}}'(\pi)\simeq\operatorname{Isom}_{\mathrm{lac}}(\widehat{\Pi}_{g,\nu},\pi)/\widehat{\unsure{\mathfrak S|\mathfrak G}}_{g,\nu}
\tag{\unsure{8}}
\]

\mspage{319}
